\section{Introduction}
\label{sec:intro}

Rolling-element bearings fail in service, and the vibration signature of a spall or a crack is well understood: an impact
train at a kinematic frequency fixed by the bearing geometry and the shaft speed, modulated by load and smeared by cage slip
\cite{randall2011}. Condition-monitoring practice turns that signature into a decision --- replace the bearing, or leave it
--- and the value of an automated diagnosis is the reliability of that decision on a machine the diagnostic system has not
seen before.

Data-driven diagnosis has moved towards transfer learning to meet exactly that requirement. Unsupervised domain adaptation
transfers a model trained under one operating condition to another without target labels \cite{zhao2021udtl}, and source-free
domain adaptation (SFDA) goes further: during adaptation only the trained source model and unlabelled target recordings are
available, which suits a fleet operator who cannot ship raw data between sites \cite{liang2020shot,sdalr2025}. Reported
accuracies are high. In a survey of 18 bearing-diagnosis papers published in 2025, every one reported above
\SI{94}{\percent} and six reported \SI{100}{\percent} \cite{vieira2026}. None of the 18 evaluated on bearings that were
absent from training.

That omission matters because a vibration dataset carries two kinds of information at once. One is the fault physics that
should transfer: the impact train at the outer-race or inner-race passing frequency. The other is the acoustic identity of an
individual specimen --- its clearance, mounting, surface finish, resonances and the particular way its damage was produced.
When a model may see the same physical bearing in training and in test, identity is sufficient to score well, and no
experiment in that protocol can tell the two apart. Splitting by physical bearing removes the shortcut, and in supervised
single-bench settings it has already been shown to change conclusions \cite{hendriks2022,wheat2024leakage,vieira2026}.

SFDA deserves separate scrutiny, for a reason specific to how it works. Adaptation is driven by the source model's own
predictions on target data: pseudo-labels, centroids or neighbourhood structure, refined over iterations. If the source
representation has learned to recognise the individual bearings it was trained on, then on a new bearing there is nothing
correct to refine --- the procedure sharpens a decision that was never about the fault. Whether this happens, how much
performance it costs, and whether anything can repair it are empirical questions that a leaky protocol cannot ask.

Physics offers a natural remedy and is increasingly used as one: characteristic-frequency evidence can filter pseudo-labels
during adaptation, as in the EAGLE module of Bio-SFDA \cite{yoon2026biosfda} or the consistency validator of PCTL
\cite{pctl2026}. Such filters are usually justified by end accuracy on the same leaky protocols. Their false acceptance on
healthy bearings --- the rate at which they certify a fault that is not there, which is what makes them safe or unsafe inside
an adaptation loop --- is rarely reported, and the amount of damage a perfect filter could repair has not been measured.

This paper follows one thread through those questions, with the evaluation apparatus fixed before the experiments were run.

\begin{enumerate}
  \item \textbf{How much survives.} A paired design holds the source model and the target operating condition fixed and
        changes only bearing overlap, so the leaky--clean difference is attributable to that overlap alone
        (Section~\ref{sec:protocol}). We measure it on two mechanical folds, ten pre-registered random bearing splits and a
        second test rig, for two SFDA methods and for non-adaptive baselines (Sections~\ref{sec:res-collapse}
        and~\ref{sec:res-second}).
  \item \textbf{What the model learned instead.} Saved per-window predictions and a linear probe on the frozen source
        features locate the failure: the representation separates same-class bearings almost perfectly, and the adapted
        model assigns one label to each physical bearing (Sections~\ref{sec:res-mechanism} and \ref{sec:res-probe}).
  \item \textbf{What could be repaired at best.} An oracle filter that keeps only correct pseudo-labels splits the collapse
        into a filter-recoverable part and a part carried by the source representation, which bounds every pseudo-label
        filter, physical or statistical (Section~\ref{sec:res-decomp}).
  \item \textbf{What physics gating actually buys.} Four gates are compared by false acceptance on 480 real healthy
        recordings and over their full operating curves, then inserted into the adaptation loop and evaluated as paired
        differences against the oracle ceiling (Section~\ref{sec:res-gating}).
  \item \textbf{Kinematic housekeeping.} We report per-manufacturer Paderborn geometry, a measured cage-slip distribution
        that settles when 6203 fault lines can lock onto shaft orders (Section~\ref{sec:res-slip}), and dataset corrections
        in the supplementary material.
\end{enumerate}

We claim no new network, no new detector and no new adaptation algorithm. Statistical envelope thresholds are established
\cite{borghesani2013,randall2011}, harmonic-coordinate features close to the ones used here were published by Matania et al.
\cite{matania2025}, and class-conditional adversarial invariance is standard in domain generalisation \cite{li2018ciddg}. The
contribution is the leakage-controlled evaluation of source-free adaptation, the measured mechanism behind its failure, the
ceiling that bounds physics gating, and the evidence a result in this area should carry before it is believed.
